\uuid{QGdU}
\exo7id{7445}
\titre{exo7 7445}
\auteur{mourougane}
\organisation{exo7}
\datecreate{2021-08-10}
\isIndication{false}
\isCorrection{false}
\chapitre{Géométrie affine dans le plan et dans l'espace}
\sousChapitre{Géométrie affine dans le plan et dans l'espace}
\cpge{mpsi-mp2i.matrices-applications-lineaires.a, pcsi.matrices-applications-lineaires.a, ptsi.matrices-applications-lineaires.a, mpsi-mp2i.espaces-prehilbertiens.d, pcsi.espaces-prehilbertiens.d, bcpst1.applications-lineaires-matrices, mp-mpi.endomorphismes-euclidien.b, pc.endomorphismes-euclidiens.b, psi.endomorphismes-euclidien.b, pt.isometries-euclidien.b, mp-mpi.endomorphismes-euclidien.e, psi.endomorphismes-euclidien.e, pt.isometries-euclidien.d, bcpst2.applications-lineaires-matrices.c}
\module{Géométrie}
\niveau{L2}
\difficulte{}

\contenu{
\texte{
Soit $(E,<.>)$ un espace vectoriel euclidien. Répondre par vrai ou
faux, puis dire sous quelles hypothèses supplémentaires sur la base
l'affirmation est vraie.
}
\begin{enumerate}
    \item \question{La matrice $A$ d'une application linéaire orthogonale $u$ dans une base $\mathcal B$ de $E$ vérifie $\mbox{}^tAA=Id$.}
    \item \question{Dans toute base orthonormée de $\Rr^3$, la matrice d'une
rotation  est  de   la  forme  $  \left(\begin{array}{ccc}
1&0&0\\0&\cos\theta&-\sin\theta\\0&\sin\theta&\cos\theta
\end{array}\right)
$}
\end{enumerate}
}
